\chapter{绪论}
\section{研究背景}

当前诸多的人工智能系统归根结底是依赖于数据中的统计相关关系。这种基于相关性建模范式的人工智能系统在各种领域无处不在，如大语言模型、推荐系统、预测模型，但其内在局限性日益凸显。仅依赖相关关系的人工智能系统存在诸多深层次问题：首先，相关性无法揭示变量间的内在作用机制，导致模型缺乏可解释性，难以回答"为什么"这类根本性问题；其次，相关关系对数据分布变化极为敏感，当训练数据与测试数据存在分布偏移时，模型性能会急剧下降，泛化能力严重受限；再者，相关性容易受到混杂因素的干扰，产生虚假关联，使得决策系统可能基于错误的逻辑链条进行推理；最后，纯相关性模型无法支持反事实推理，限制了其在需要因果推断的复杂决策场景中的应用潜力。

因果关系作为解决上述问题的核心概念，其本质在于揭示变量间的定向作用机制，即一个变量的变化如何引起另一个变量的变化，这种作用机制独立于观测分布而存在，具有方向性、不变性与可干预性三大特征 \cite{Pearl2009, Peters2017}。因果关系不仅关注"是什么"，更深入探究"为什么"和"会怎样"，为人工智能系统提供了理解世界运行规律的新范式。
因果关系在智能决策、系统鲁棒性和知识发现方面具有深远意义。在现实世界应用中，其重要性主要体现在以下三个方面：
第一，增强决策的可靠性与可解释性。因果模型能够识别真正的驱动因素，避免被表面相关性误导，使人工智能系统在医疗诊断、金融风控等高风险领域做出更加可靠且可解释的决策，满足监管要求和用户信任需求。
第二，提升系统的适应性与泛化能力。基于因果机制的模型能够识别数据中的不变规律，即使在环境变化或数据分布偏移的情况下，仍能保持稳定的性能表现，这对于自动驾驶、工业控制等需要长期稳定运行的系统至关重要。
第三，支持主动干预与策略优化。因果推理使人工智能系统能够预测不同干预措施的效果，支持反事实推理，为个性化推荐、精准营销、政策制定等需要主动干预的场景提供科学依据，实现从被动预测到主动优化的转变。

因果关系可以分为因果发现、因果推断与反事实推理三大主流研究分支，这一体系与Pearl提出的因果关系之梯高度一致，分别对应观察、干预与想象三个递进层次 \cite{Pearl2009}。因果发现致力于从观测或干预数据中自动识别变量间的因果结构，包括因果对的判定、因果方向的确定以及完整因果图的构建，典型应用涵盖基因调控网络构建与复杂工业系统的故障溯源\cite{Cavenaghi2024}；因果推断则建立在已知因果结构的基础上，重点在于量化原因变量变化对结果变量的影响程度，广泛应用于药物疗效评估与公共政策效果分析等领域 \cite{Garrido2021}；反事实推理作为因果推理的最高层次，旨在探讨与事实相反的情境下结果可能发生的变化，该方法要求模型能够刻画潜在结果或基于结构因果模型进行完备建模，典型应用包括法律责任判定、个性化医疗决策支持以及算法公平性审计等 \cite{Parafita2022}。因果发现产生了因果结构和函数，是因果推断和反事实推理的基石，本文研究聚焦于因果发现。

{\Huge 这一段还要重写}因果发现旨在从数据中精准挖掘变量间的真实因果关联，厘清因果作用的方向与路径，构建能够客观反映变量内在依赖机制的因果结构模型。发现因果结构模型为后续因果效应量化、反事实推理等任务提供可靠的结构基础，并且对推动人工智能从“关联驱动”向“因果驱动”升级具有至关重要的意义 \cite{Squires2023, Yao2021}。当前，因果发现方法存在两种主流的分类范式：第一种依据方法的核心思想将因果发现方法划分为基于约束的方法、基于评分的方法以及基于结构不对称性的方法；第二种则从结构识别的实现逻辑出发，将因果发现方法分为两大类，即通过组合和搜索算法识别图形结构的方法，以及基于连续优化识别图形结构的方法。本文的研究重点聚焦于第二种分类范式中的后者，即基于连续优化的图形结构识别方法。从实现逻辑与性能特性来看，这两类结构识别方法存在显著差异：通过组合和搜索算法识别图形结构的方法，核心思路是先生成大量候选因果图形结构，再通过独立性检验或评分函数筛选出最优结构，其本质是在离散的图形结构空间中进行搜索；这类方法的优势在于理论逻辑直观、可解释性较强，但随着变量维度的提升，候选图形结构的数量呈指数级增长，导致搜索空间急剧扩大，不仅计算复杂度激增、效率低下，还极易陷入局部最优解，难以适用于高维变量场景。与之相对，基于连续优化识别图形结构的方法，核心思路是将离散的因果图形结构（通常用邻接矩阵表征）转化为连续的优化变量，通过构建含无环性等约束条件的连续可微目标函数，利用梯度下降等高效优化算法求解最优解，进而得到对应的因果结构 \cite{Zheng2020}；该类方法有效突破了组合和搜索算法在高维场景下的瓶颈 \cite{Chowdhury2024}，具备三大核心优势：一是将离散搜索转化为连续优化，大幅降低了高维场景下的计算复杂度，提升了结构识别的效率\cite{Zheng2020, Chowdhury2024}；二是可灵活融入正则化机制，便于实现结构稀疏性约束，更符合真实世界中因果结构“稀疏连接”的特性 \cite{Rashid2022, Fang2023}；三是易于与神经网络等数据驱动模型深度融合，能够更好地捕捉变量间的复杂非线性因果关联，显著提升复杂数据场景下的结构识别精度，因此成为近年来高维因果发现领域的研究热点\cite{Rashid2022, Fan2022, Deng2025}。

然而，传统的因果发现方法，特别是基于连续优化的因果发现方法主要聚焦于发现一个有向无环图（Directed Acyclic Graph, DAG），该有向无环图与特征内涵的因果关系相容，即有向无环图表示了特征之间的因果关联 \cite{HeinzeDeml2018}，但是这种依赖于有向无环图的因果关系表示方法仅能识别特征之间是否存在因果关联，无法揭示因果特征之间的内在机理，也就是因果函数（Causal Function, CF）\cite{Greenland2002}。事实上，图模型和结构方程模型是两种不同的表示因果关系的方法，前者主要使用有向无环图刻画因果，而后者使用结构方程的形式 。图\ref{有向无环图与函数因果模型之间的关系}展示了有向无环图与结构方程模型之间的关系，有时，结构方程模型也被称为函数因果模型（Functional Causal Model, FCM），顾名思义，FCM肯定了CF在因果模型中的作用，CF是因果关系的重要组成部分 \cite{Bongers2021}。对比可以发现，有向无环图缺失了CF的表示，以结构汉明距离（表现为缺边、多边和反向边三类边的和）为核心评价标准的传统方法在目标上有所欠缺。加入CF具备多方面优势：一是能够精准刻画因果作用机理，明确原因变量对结果变量的具体影响模式（如线性影响、非线性影响），提升因果模型的解释性；二是增强因果模型的干预与预测能力，基于因果函数可直接计算干预原因变量后的结果变化，为精准决策提供量化依据；三是提升模型的泛化性能，因果函数刻画的是变量间的本质依赖关系，相较于仅识别结构的DAG模型，在数据分布发生变化的场景下仍能保持较好的性能稳定性；四是为反事实推理提供支撑，通过因果函数可重构“若原因变量处于其他状态”的反事实场景，进而评估不同决策方案的潜在效果 \cite{Pearl2009}。
\begin{figure}[H]
	\centering
	\includegraphics[width=\linewidth]{有向无环图与函数因果模型之间的关系}
	\caption{有向无环图与函数因果模型之间的关系}
	\label{有向无环图与函数因果模型之间的关系}
\end{figure}

在上述动机的驱使下，本文主要探究以函数因果模型为目标的因果发现。具体而言，本文引入了一种可解释性强且符号化的神经网络Kolmogorov-Arnold Networks（KAN），并将其作为求解器融入到基于连续优化的因果发现方法之中，通过Kolmogorov-Arnold表示定理来表示因果函数。KAN的核心原理基于Kolmogorov-Arnold表示定理，该定理指出：任意多元连续函数都可分解为单变量函数的叠加与复合形式 \cite{Ziming2025}。KAN正是基于这一原理构建的神经网络模型，其网络结构由输入层、隐藏层与输出层组成，隐藏层通过一系列单变量基函数对输入变量进行映射，输出层则对隐藏层的映射结果进行线性组合，最终实现对多元函数的精准拟合；同时，KAN的结构具备天然的符号化特性，可通过稀疏正则化等方法筛选出关键的基函数与组合关系，从而实现因果函数的可解释性表达。本文致力于从观测数据中获得与特征分布相容的FCM，为因果发现领域提供一种面向可解释性和干预的研究思路。

本文主要面向真实世界应用场景的需求，研究基于KAN的因果发现的理论与方法。本章接下来的内容将会围绕这一主题展开：1.2节首先回顾国内外在因果发现领域和Kolmogorov-Arnold表示定理驱动下的神经网络研究现状的研究工作，1.3节和1.4节分析该领域目前有待解决的问题以及本文提供的解决思路，最后1.5节给出本文的组织结构并简介每章的具体内容。


\section{国内外研究现状}

本文的国内外研究现状由因果发现和Kolmogorov-Arnold表示定理驱动下的神经网络两大部分构成。1.2.1节首先介绍基于观测数据的因果发现方法，并介绍基于观测数据的方法分类标准以及每一类算法的技术特点和设计思想。1.2.2节介绍其中一种基于先验知识融合的因果发现方法，并介绍该方法的总体框架和核心思想。1.2.3节对当前研究现状进行分析与总结，分析领域内的空白并接引本文的研究内容。

\subsection{基于观测数据的因果发现方法}

因果发现作为探究因果关系的重要一环，其核心价值在于从观测数据中剥离虚假关联、还原变量间的真实因果结构，为后续因果推断、决策优化提供可靠的结构基础，是连接数据驱动与因果推理的核心桥梁。因果发现的内涵是通过数据驱动的方法识别变量间的因果依赖关系并构建因果图模型，意义在于打破传统统计学习中“关联不等于因果”的局限，使模型能够从“预测是什么”向“解释为什么”进阶，进而支撑更可靠的干预决策与跨场景泛化。在因果领域内，因果发现有两种分类方法，其一是依据方法论可以分为基于约束的方法、基于评分的方法、利用结构不对称的方法、混合方法四类 \cite{Vowels2022}；其二是依据数据空间类型可以分为基于组合和搜索的因果发现方法和基于连续优化的方法。两种分类方法互有交叉，本文采用第二种分类方法。

{\hei 1.基于组合和搜索的因果发现方法}

表~\ref{causal_discovery_methods}列举了基于组合和搜索的因果发现方法，该类方法包含了基于约束的方法、基于评分的方法、利用结构不对称的方法、混合方法四类。此外，表还提供了每种算法与相关的假设的关系，包括充分性（即是否假设不存在隐藏变量）、忠实度（某些方法实现较轻或宽松的忠实度），以及无循环性（一些方法可以学习反馈环和循环）。最后显示了方法的输出图类型（DAG、CPDAG、PAG等。）。

基于约束的因果发现方法的核心思想是通过检验数据中的条件独立性约束来构建因果图结构。其基本逻辑为：先假设变量间存在全连接的初始结构，再通过一系列条件独立性检验逐步删除不满足约束的边，最终得到与数据中条件独立性一致的因果图或马尔可夫等价类 \cite{Glymour2019}。
该方法不预设具体的函数模型形式，对数据分布的假设相对宽松，尤其在高维变量场景下展现出良好的可扩展性\cite{Schoelkopf2021}。然而，其有效性高度依赖于条件独立性检验的可靠性，在样本量有限时易产生偏差，且在处理非线性关系及复杂图结构时面临挑战\cite{Schoelkopf2021}。
代表性方法包括PC算法、FCI算法等。其中PC算法通过逐步筛选变量的条件集来检验条件独立性，进而确定边的存在性与方向\cite{Spirtes2000}；
FCI算法则在PC算法基础上拓展，能够处理存在隐藏变量与选择偏倚的场景，输出部分祖先图（PAG）\cite{Spirtes1995}；
在PC和FCI的基础之上，后续开发了一系列的改进算法。
KCI\cite{Sun2007}与KCI-test\cite{Zhang2012}从检验方法泛化性切入，RFCI\cite{Colombo2011}优化隐变量处理逻辑，Parallel-PC\cite{Le2019}提升高维计算效率，PC-stable\cite{Colombo2014}增强算法稳定性，FCI-soft\cite{Kocaoglu2020}与psi-FCI\cite{Jaber2020}则分别改进约束阈值与偏倚处理机制。
基于约束的因果发现方法的优势在于无需预设具体的因果模型形式，对数据分布的假设较为宽松，且在高维变量场景下具有一定的可扩展性；局限在于过度依赖条件独立性检验的可靠性，当样本量较小时检验结果易出现偏差，进而导致图结构推断错误，同时难以处理变量间的非线性依赖关系，且在复杂图结构中对条件集的选择较为敏感\cite{Vowels2022}。

基于评分的因果发现方法将因果发现转化为图结构的搜索与优化问题，核心是定义一个评分函数来量化候选因果图与观测数据的契合程度，再通过搜索算法找到评分最优的图结构。其核心框架包括两部分：一是评分函数的设计，二是高效的搜索策略。
常用的评分函数有贝叶斯信息准则（BIC）\cite{Geiger1994}、贝叶斯 狄利克雷等价（BDe）评分\cite{Heckerman1995}、最小描述长度（MDL）\cite{Adriaans2008}等，其中BIC评分兼顾模型拟合度与复杂度，可有效避免过拟合；BDe评分适用于离散数据，需预设先验分布。
搜索策略则用于解决候选图结构的组合爆炸问题，常用的有贪婪搜索\cite{Cooper1992}、爬山法\cite{Heckerman1995}、遗传算法\cite{Elidan2008}等，
进阶方法如贪婪等价搜索（GES）\cite{Chickering2003}则通过搜索马尔可夫等价类来提升效率；
在GES的基础之上，开发了许多改进算法，包括BC\cite{ElGhaoui2008}与TC\cite{Pellet2008}优化评分函数的正则化机制，Adaptive Lasso\cite{Michailidis2010}实现因果结构的稀疏性约束，GIES\cite{Hauser2012}拓展隐变量场景下的等价类搜索，Pen-PC\cite{Ha2015}融合评分与约束机制提升推断精度，K-A*\cite{Scanagatta2016}创新启发式搜索策略加速高维图优化，GES-mod\cite{AlonsoBarba2013}改进等价类方向判定规则；值得一提的是，CAM\cite{Buehlmann2013}算法通过将加性因果机制嵌入评分函数设计，并利用加性噪声模型的非对称性实现因果方向的直接识别，同时以结构约束缩小搜索空间。
基于评分的因果发现方法的优势在于直观易懂，能够直接量化图结构与数据的匹配程度，且对条件独立性检验的依赖较低，可处理部分非线性场景；局限在于搜索过程易陷入局部最优解，尤其在变量数量较多时，搜索效率大幅下降，同时评分函数的选择对结果影响较大，且难以确定因果图的方向（通常输出马尔可夫等价类，如CPDAG）\cite{Vowels2022}。

利用结构不对称的方法的核心思路是利用因果关系中存在的结构不对称性来识别因果方向，突破了传统方法难以区分马尔可夫等价类中因果方向的局限。其基本逻辑为：在因果关系X→Y中，原因X与结果Y的生成机制存在不对称性（如噪声分布、函数复杂度、边际分布与条件分布的独立性等），这种不对称性可通过数据挖掘并用于因果方向推断。代表性方法包括基于非高斯加性噪声模型（LiNGAM）\cite{Shimizu2006}、后非线性加性噪声模型（Non-linear ANM）\cite{Patrik2008}、信息几何因果推断（IGCI）\cite{Janzing2012}等，其中LiNGAM假设数据生成过程为线性且噪声非高斯，通过独立成分分析（ICA）分离噪声，进而确定因果方向；Non-linear ANM凭借非线性因果函数与加性噪声的结构不对称设计，实现了从线性到非线性因果场景的方向推断拓展；IGCI则利用“原因的边际分布与因果机制独立”的假设，通过信息几何度量来区分因果方向；以及PNL算法\cite{Zhang2009}凭借 “线性变换 + 非线性映射 + 加性噪声” 的复合结构设计，拓展了结构不对称方法在复杂非线性因果中的适用性。利用结构不对称的方法的优势在于无需干预数据即可实现因果方向的识别，为解决“因果推断基本问题”提供了新思路；局限在于需要较强的模型假设，这些假设在真实场景中可能不成立，且大多为局部方法，难以直接构建大规模因果图。

除了以上三类方法之外，MMHC\cite{Tsamardinos2006}、ARGES\cite{Nandy2016}等算法融合了基于约束和基于评分方法的核心思想，作为一种兼顾两种思想优势的混合方法。总体而言，基于组合和搜索的因果发现方法将图形结构DAG视为离散组合对象，图形的核心是 “节点间有无边、边的方向” 等离散属性，学习目标是从 “所有可能的离散结构集合”中筛选出符合数据特征的结构。可能的DAG数量随着变量数量的增加呈超指数增长\cite{Harary1973}，正如Peters等人\cite{Peters2017}指出，10个变量的可能DAG数量$>4\times10^{18}$，该搜索问题是NP-hard问题\cite{Chickering1996}。因此，传统因果发现方法在低维数据场景下具备良好的性能和可解释性，但普遍面临高维数据下搜索空间爆炸、强假设依赖与真实场景脱节、难以揭示因果内在机制等问题。
\begin{table}[H]
	\centering
	\begin{minipage}[t]{\linewidth}
		\caption{基于组合和搜索的因果发现算法}
		\label{causal_discovery_methods}
		\begin{tabularx}{\linewidth}{lcccccc}
			\toprule[1.5pt]
			\textbf{算法} & \textbf{年份} & \textbf{类型} & \textbf{充分性} & \textbf{忠诚性} & \textbf{无循环性} & \textbf{输出类型} \\
			\midrule[1pt]
			PC\cite{Spirtes2000} & 1993 & 约束 & 是 & 是 & 是 & CPDAG \\
			FCI\cite{Spirtes1995} & 1995 & 约束 & 否 & 是 & 是 & POIPG \\
			KCI\cite{Sun2007} & 2007 & 约束 & 是 & 是 & 是 & CPDAG \\
			KCI-test\cite{Zhang2012} & 2012 & 约束 & 是 & 是 & 是 & CPDAG \\
			RFCI\cite{Colombo2011} & 2012 & 约束 & 否 & 是 & 是 & PAG \\
			Parallel-PC\cite{Le2019} & 2014 & 约束 & 是 & 是 & 是 & CPDAG \\
			PC-stable\cite{Colombo2014} & 2014 & 约束 & 是/否 & 是 & 是/否 & CPDAG \\
			FCI-soft\cite{Kocaoglu2020} & 2019 & 约束 & 否 & 松弛 & 是 & MEC \\
			psi-FCI\cite{Jaber2020} & 2020 & 约束 & 否 & 松弛 & 是 & Psi-EC \\

			EG\cite{Elidan2008} & 2009 & 评分 & 是 & 是 & 是 & BT-DAG \\
			CAM \cite{Buehlmann2013} & 2014 & 评分 & 是 & 是 & 是 & CPDAG \\
			K2\cite{Cooper1992} & 1992 & 评分 & 否 & 是 & 是 & CPDAG \\
			HGC\cite{Heckerman1995} & 1995 & 评分 & 是 & 是 & 是 & CPDAG \\
			GES\cite{Chickering2003} & 2002 & 评分 & 是 & 是 & 是 & CPDAG \\
			BC\cite{ElGhaoui2008} & 2008 & 评分 & 是 & -- & 否 & UG \\
			TC\cite{Pellet2008} & 2008 & 评分 & 是 & 是 & 是 & CPDAG \\
			Adaptive Lasso\cite{Michailidis2010} & 2010 & 评分 & 是 & 是 & 是 & DAG \\
			GIES\cite{Hauser2012} & 2012 & 评分 & 是 & 是 & 是 & PDAG \\
			GES-mod\cite{AlonsoBarba2013} & 2013 & 评分 & 是 & 是 & 是 & CPDAG \\
			Pen-PC\cite{Ha2015} & 2015 & 评分 & 是 & 是 & 是 & CPDAG \\
			K-A*\cite{Scanagatta2016} & 2016 & 评分 & 是 & 是 & 是 & DAG \\

			LiNGAM\cite{Patrik2008} & 2006 & 结构不对称 & 是 & 否 & 是 & DAG \\
			non-linear ANM \cite{Janzing2012} & 2008 & 结构不对称 & 是 & 是 & 是 & DAG \\
			PNL {Zhang2009} & 2009 & 结构不对称 & 是 & 否 & 是 & DAG \\
			MMHC \cite{Tsamardinos2006} & 2006 & 混合 & 是 & 是 & 是 & DAG \\
			ARGES \cite{Nandy2016} & 2018 & 混合 & 是 & 是 & 是 & CPDAG \\

			\bottomrule[1.5pt]
		\end{tabularx}
	\end{minipage}
\end{table}


基于连续优化的因果发现方法主要属于基于评分的因果发现方法的范畴，近年来也有采用约束或利用不对称的连续优化方法。不同于基于组合和搜索的因果发现方法，该因果发现方法是应对高维、复杂数据场景的核心技术创新，其核心思路是将将离散结构学习转化为连续优化问题，通过设计连续约束与损失函数，利用梯度下降、增强拉格朗日法等优化工具求解最优因果结构，彻底规避了传统组合和搜索方法的NP-hard困境。 近年来，基于连续优化的因果发现方法被广泛研究，研究人员不仅继续完善基于连续优化的因果发现方法理论，同时将算法拓展到现实世界中的各类因果发现挑战中，取得了丰硕的成果。

连续优化方法在深度学习领域非常普遍，通过梯度下降的变体来优化高参数化的网络\cite{Goodfellow2016}。神经网络的动机在于它们不先验限制功能形式，因此“让数据自己说话”\cite{VanderLaan2011}。计算能力的提升（尤其是GPU的出现）使得从大型高维数据集中学习变得可行。近年来，越来越多的方法试图从数据中学习结构，同时利用持续优化的优势，这促成了黑箱深度学习方法与因果发现的融合。基于连续优化的方法可以分为具有无环约束的方法，缺乏无环约束的方法，以及其他类型的方法。

具有无环约束的方法是连续优化方法的主流技术分支，其核心目标是在连续优化框架内严格保证学习到的因果结构为有向无环图，这也是因果发现中最核心的结构要求之一。该分支的技术奠基始于Zheng等人2018年提出的NO TEARS算法\cite{Zheng2018}，该算法被广泛认为是首个成功将离散组合图搜索问题重新定义为连续优化问题的开创性工作，彻底突破了传统组合搜索方法面临的NP-hard困境。其核心实现思路是通过矩阵迹运算将离散的无环性约束转化为可微分的连续条件，通过增强拉格朗日法求解带约束的优化问题。但NO TEARS尚存在以下三个不足之处：一是矩阵指数运算的计算复杂度高达$O(d^3)$，难以适配高维数据场景；二是无环性约束采用连续近似，实际求解中$h(A)+$
可能非零，需通过阈值化处理确定边的存在，易引入结构偏差；三是线性求解器仅能建模线性因果关系，适用场景受限。

NO TEARS提出的无环约束核心思想为后续研究奠定了基础，具有无环约束的方法围绕“计算效率优化”“突破线性限制”“提升约束精度”三大方向形成清晰发展脉络。在计算效率优化方面，NO BEARS算法\cite{Lee2019}采用谱半径近似策略替代NO TEARS中的矩阵指数运算，在完整保留连续无环约束核心逻辑的前提下，将计算复杂度从$O(d^3)$显著降至$O(d^2)$，为高维数据因果发现提供了效率支撑；CASTLE算法\cite{Kyono2020}通过正则化辅助的连续优化，引入辅助因果图引导主结构学习，专门适配异质性/非平稳高维数据场景，提升计算效率；LEAST算法\cite{Zhu2021}则在NO BEARS的基础上进一步深化，推导谱半径上界表达式，通过更严谨的理论转化将无环性约束的计算复杂度从$O(d^2)$降至线性级别$O(d)$，能够适配16万节点及以上的超大规模数据场景；DAGOR\cite{Zuo2024}引入了图拓扑排序来解决连续优化问题，该问题远小于DAG空间，有助于避免局部最优且显著降低计算成本，并提出了基于QR分解的新DAGs学习优化策略。

在突破线性限制方面，DAG-GNN算法\cite{Yu2019}通过整合神经网络函数$f$和黑箱变分推断，扩展了NO TEARS，通过VAE学习变量低维表征实现非线性因果关系的隐式建模。同年提出的GAE算法\cite{Ng2019}承接这一思路，采用图自编码器架构，将连续无环约束融入编码器-解码器的联合优化过程，通过图结构重构误差与无环约束的双目标优化提升结构学习精准性，丰富了表征转换路径的技术选择。NO TEARS+算法\cite{Zheng2020a}则开辟了直接非线性建模路径，在NO TEARS原有的连续优化框架基础上，引入非参数化模型（如核方法、神经网络）替代线性结构方程模型，无需依赖低维表征转换即可直接建模复杂非线性因果关系，简化了非线性建模流程。此外，CausalVAE算法\cite{Yang2021}作为表征转换路径向非语义数据场景的延伸，引入因果层与语义监督，将连续优化得到的因果结构与高维非语义数据（如图像）表征相结合，使因果结构具备清晰的语义关联，进一步拓展了突破线性限制方向的应用范围。

在提升约束精度方面，GOLEM算法\cite{Ng2020}针对NO TEARS中稀疏性与无环性约束协同不足的问题，设计加权损失函数优化两者的权重分配，平衡了模型拟合度、稀疏性与无环性三者的关系，有效减少边误判，成为线性场景高精度学习的代表性算法；同期的NO FEARS算法\cite{Wei2020}则聚焦约束本身的近似精度，设计更精细的连续近似函数优化无环性约束的数学表达，降低了连续变量与离散无环性之间的映射偏差，进一步提升了结构识别的可靠性；DAGMA算法\cite{Bello2022}将对数行列式和M举矩阵概念引入到无环性约束上构建无环约束，并证明约束的所有驻点均为全局极小值，且这些驻点对应于有向无环图；LeCaSiM\cite{Cai2024}给出DAG及其对应简要证明的一系列代数等价条件，然后利用M矩阵的性质推导出新的无环约束函数，用于克服优化过程容易出现数值爆炸或高阶约束条件梯度消失，导致收敛条件不稳定，结果偏离真实结构等问题。

除上述主流方向外。GranDAG算法\cite{Lachapelle2020}则采用非线性加法噪声结构模型，形式为$X_j = f_j \cdot pa_j+ U_j$，其中每个函数$ f_j $参数化为全连通神经网络。为了保持机制独立性，这与邻接矩阵所暗示的独立性相对应，他们提出了神经网络路径和连通矩阵，类似于Germain等人\cite{Germain2015}的早期工作。SAM算法\cite{Kalainathan2022}通过使用对抗训练解决了这两个局限性，并将优化DAG的机制纳入端到端训练。他们的得分函数是一个带有两个模型复杂度正则化器的对数似然损失：一种是按顶点对模型进行与顶点父数成正比的惩罚;以及一种作为神经网络参数/权重衰减的功能。MaskedNN算法\cite{Ng2022}则还尝试利用神经网络改进NO TEARS，他们同样假设加法噪声SEM，并解释了他们的方法如何从处理标量变量直接扩展到向量值变量。它们讨论了可识别性（这是组合和连续优化文献中许多方法常被省略的），他们还强调，DAG-GNN和GAE中邻接矩阵的非线性变换可能影响其因果可解释性。针对时序数据的特殊需求，DYNOTEARS算法\cite{Pamfil2020}将NO TEARS的连续无环约束框架拓展至时序场景，设计双邻接矩阵，通过对双邻接矩阵施加无环约束，实现时序数据中同期效应与滞后效应的同步捕捉，成为连续优化因果发现向时序场景拓展的核心成果；Dagma-DCE算法\cite{Waxman2024}理论证明“独立性”的不透明代理可能与实际因果强度存在任意差异，使用可解释的因果强度度量来定义加权邻接矩阵，增强了连续可微因果发现的可解释性Jeroen等人\cite{Berrevoets2023}认为基于连续优化的方法在承认可微性的同时也带来了可迁移性，并提出了D\mbox{-}Struct，通过一种新颖的架构和损失函数恢复了发现结构中的可迁移性，同时保持完全可微。

缺乏无环约束的方法是基于连续优化的另一种技术分支。因果马赛克\cite{Wu2020}是一种基于神经网络的非线性独立成分分析，旨在区分双变量对之间的因果方向。同样，AEQ方法\cite{Galanti2020}的作者基于自编码器的重建误差开发了一个评分函数，用于发现向量值因果对的方向性，他们通过基于原始切片的排序连接创建变量的多变量版本，将该结果扩展到单变量$X$。该多变量代理的复杂度随后通过自编码器重建误差（使用$L_2$损耗）来测量。CGNN\cite{Goudet2018}将神经网络与爬山法或禁忌搜索结合起来，他们的方法包括识别可能的隐藏混杂因素的方法，利用混杂现象可以被建模为（否则）潜在潜在随机变量之间的相关性/关联。DEAR\cite{Shen2020}，它结合了变分自编码器（VAE）\cite{Kingma2013,Rezende2014}和对抗损耗\cite{Goodfellow2014}，以推断具有“因果”结构的潜在空间，假设已经给出了“超级图”，并学习了相关的权重和参。CAREFL\cite{Khemakhem2021}将因果发现与称为归一化流的深度学习框架结合，归一化流提供了构建生成模型的方法，这些模型能够通过基本且可解密度的可逆变换来建模复杂密度。它们通过变量变更公式和逆对数雅可比行列式，精确计算构成其学习目标的对数似然。ICL\cite{Wang2020}的作者聚焦于缺失数据环境下的结构发现问题，利用不完整的数据同时通过GAN补入缺失数据，以匹配生成的分布与经验分布。SDI\cite{Ke2019}假设忠实性和充分性，且限制在离散的类别变量中，且无缺失，利用经历未知干预的数据来发现结构。Ke等人\cite{Ke2020}提出了一种元学习神经网络方法，该方法还包含干预并利用图的连续表示，能够快速且高效地学习新的因果模型。CAN\cite{Moraffah2020}在推理时便于从因果图进行介入采样，在图像数据和传统数据上展现竞争性能。RL-BIC的作者\cite{Zhu2019}明确采用强化学习方法进行因果发现。它们使用编码\mbox{-}解码神经网络模型生成有向图。编码\mbox{-}解码器的输出是拟议的图，利用BIC进行评分以生成奖励信号。

以及其他类型的方法指的是既不使用无环约束，也与神经网络无关的方法。Vrando（2020）\cite{Varando2020}提出了一个近端梯度优化目标，能够得到线性SEM和相应的DAG。该方法通过将学习问题框架为稀疏矩阵分解来推导新颖的目标，最终的NODAG方法被证明既有效又高效。ABIC\cite{Bhattacharya2021}扩展了连续优化范式，发现了各种类型的 ADMG（祖先型、干旱型、无弓型）来解释未测量的混杂因素，并提出了三个可微约束，可用于发现特定类型的ADMG。表\ref{tab:Continuous_causal_methods}总结了具有代表性的基于连续优化的因果发现方法，包括算法面向的数据特点，使用的求解器模型，是否遵循无循环性和最后的输出类型。值得一提的是，表中的无循环性指的是算法输出的结果是否保证无循环性，而不是是否使用了无循环约束。
\begin{table}[H]
	\centering
	\begin{minipage}[t]{\linewidth}
		\caption{基于连续优化的因果发现方法}
		\label{tab:Continuous_causal_methods}
		\begin{tabularx}{\linewidth}{lcccccc}
			\toprule[1.5pt]
			\textbf{算法} & \textbf{时间} & \textbf{数据特点} & \textbf{模型} & \textbf{无循环性} & \textbf{输出类型} \\
			\midrule[1pt]
			NO TEARS\cite{Zheng2018} & 2018 & 低维 & 线性 & 是 & DAG \\
			CGNN\cite{Goudet2018} & 2018 & 低维 & 神经网络 & 是 & DAG \\
			SAM\cite{Kalainathan2022} & 2019 & 低维/中维 & 神经网络 & 是 & DAG \\
			DAG-GNN\cite{Yu2019} & 2019 & 低维 & 神经网络 & 是 & DAG \\
			GAE\cite{Ng2019} & 2019 & 低维 & 神经网络 & 是 & DAG \\
			NO BEARS\cite{Lee2019} & 2019 & 低维/中维/高维 & 多项式 & 是 & DAG \\
			DEAR\cite{Shen2020} & 2020 & 图像 & 神经网络 & 是 & - \\
			CAN\cite{Moraffah2020} & 2020 & 低维/中维/图像 & 神经网络 & 是 & DAG \\
			NO FEARS\cite{Wei2020} & 2020 & 低维 & 线性 & 是 & DAG \\
			DAGMA\cite{Bello2022} & 2022 & 低维 & 神经网络 & 是 & DAG \\
			Dagma-DCE\cite{Waxman2024} & 2024 & 低维 & 神经网络 & 是 & DAG \\
			DTSL\cite{Berrevoets2023} & 2023 & 低维 & 神经网络 & 是 & DAG \\
			LeCaSiM\cite{Cai2024} & 2024 & 低维 & 神经网络 & 是 & DAG \\
			DAGOR\cite{Zuo2024} & 2024 & 低维/中维/高维 & 神经网络 & 是 & DAG \\
			GOLEM\cite{Ng2020} & 2020 & 低维 & 神经网络 & 是 & DAG \\
			ABIC\cite{Bhattacharya2021} & 2020 & 低维 & 线性 & 是 & ADMG/PAG \\
			DYNOTEARS\cite{Pamfil2020} & 2020 & 低维 & 线性 & 是 & SVAR \\
			SDI\cite{Ke2019} & 2020 & 低维 & 神经网络 & 是 & DAG \\
			AEQ\cite{Galanti2020} & 2020 & 二元 & 神经网络 & - & direction \\
			CRN\cite{Ke2020} & 2020 & 低维 & 神经网络 & 是 & DAG \\
			CASTLE\cite{Kyono2020} & 2020 & 低维/中维 & 神经网络 & 是 & DAG \\
			GranDAG\cite{Lachapelle2020} & 2020 & 低维 & 神经网络 & 是 & DAG \\
			MaskedNN\cite{Ng2022} & 2020 & 低维 & 神经网络 & 是 & DAG \\
			CausalVAE\cite{Yang2021} & 2020 & 图像 & 神经网络 & 是 & DAG \\
			CAREFL\cite{Khemakhem2021} & 2020 & 低维 & 神经网络 & 是 & DAG/direction \\
			Varando\cite{Varando2020} & 2020 & 低维 & 线性 & 是 & DAG \\
			NO TEARS+\cite{Zheng2020a} & 2020 & 低维 & 非线性 & 是 & DAG \\
			ICL\cite{Wang2020} & 2020 & 低维 & 神经网络 & 是 & DAG \\
			LEAST\cite{Zhu2021} & 2020 & 低维/中维/高维 & 线性 & 是 & DAG \\
			CausalMosaic\cite{Wu2020} & 2020 & 二元 & 神经网络 & - & direction \\
			\bottomrule[1.5pt]
		\end{tabularx}
	\end{minipage}
\end{table}

\subsection{基于知识融合的因果发现方法}
纯粹基于观测数据的因果发现能力在一定程度上已触及瓶颈，其精确度极大程度地受到数据质量、样本规模及噪声分布的影响。在某些特定领域应用中，训练数据往往有限且存在较强的观测噪声，仅利用训练数据学习因果结构不仅面临可识别性理论上的困难，在实际操作中也难以保证模型的鲁棒性。先验知识的引入可以有效提高学习算法的效率与准确性，尤其是在结构因果模型（SCM）的构建过程中，因果知识能够帮助指导模型的建立过程，解决数据不充分领域的因果建模问题。此外，因果图模型提供了变量间依赖关系的图形化表示，易于被人类理解，这一关键属性开启了学习过程中人为干预的可能性。目前主流的知识融合方式主要从结构约束与机制约束两个维度利用先验知识辅助因果发现，一方面是通过建立合理可靠的先验结构引导后续的学习方向，另一方面是通过限制优化空间或函数形式来提高结构学习的效率与可解释性。

{\hei 1.结构先验的构建与引导}

因果发现算法的性能很大程度上取决于先验结构的质量，不准确的先验结构可能会误导因果预测算法，使算法产生错误的解。通过领域知识构建因果网络的先验结构，排除特定的错误边或确认存在的因果边，可以引导后续的学习过程向正确的方向前进。现有的基于先验知识的因果发现方法中，多数学者对于结构先验的选择关注较少，简单的结构先验往往是基于领域实用的需要，且往往是结构上一致的先验\cite{Meng2022}。通过领域知识定义变量之间的特定关系或者顺序结构，可在一定程度上避免依赖数据搜索过程中的错误边。有研究通过指定领域知识来确定部分存在关联和不存在关联的先验结构，减少后续搜索过程可能会出现的模型集合\cite{Hoejsgaard1995}。另有研究认为成功的决策取决于推断潜在的或潜在的信息，通过使用专家知识和可用数据的组合来预测和推理潜变量，依靠领域专长来识别重要的潜变量，并建立它们与观测变量之间的关系\cite{Yet2013}。通过一系列专家评审，将领域专家知识与可用数据进行整合，系统地开发和完善模型实用化的方法来开发带有隐变量的因果模型。

完善的领域知识不仅可以提高模型构建的准确度，还提高了模型构建的效率，节省了大量的数据和算力。然而大多数领域内的领域知识并不全面，只涵盖部分边的信息，这些信息称为部分先验信息。部分研究考虑到部分先验信息的不准确性以及高质量的领域知识难以获取，提出设置领域知识的置信度并充分利用节点之间的联合分布概率，扩展了多类领域知识，包括父节点存在、父节点不存在，以及分布知识，将领域知识及其置信度值与训练数据融合，计算建议的接受概率，学习到更好的因果网络结构\cite{Xu2015}。在一些特殊领域内，领域知识包含了模型构建的关键信息，且很大程度上可以避免风险数据造成的影响，帮助模型更好的应用于实际问题的决策推理中。例如，将医疗领域内生物信号通路中包含的知识转化为先验概率，提出了一个包含任何先验信息来源的通用框架\cite{Boluki2017}；或将物理领域内的关键知识引入，初步建立起强制关系获取先验知识模型，并在此基础上通过数据驱动学习因果结构，通过实验证明了先验信息可以防止已识别的关键原因并阻止其传播来缓解原因耦合\cite{Meng2022}。

{\hei 2.优化空间的约束与限制}

由于因果图的候选结构搜索空间过大，尤其是随着变量维度的增加呈超指数级增长，很多学者利用领域知识对搜索空间进行限制来提高搜索效率，包括限制网络节点的出入度、节点排序、节点之间的弧和边的存在性。早期研究试图解决因果网络的精化问题，即根据新的数据更新因果网络的参数和结构，通过引入专家知识，要求专家对变量进行全排序的前提下进行因果网络结构的学习。通过排序的方式使得模型搜索空间缩小了一半以上，提高了搜索的效率和准确度。另有研究利用专家知识启发进行结构学习，建立了一种概率框架对专家两两比较数据进行建模\cite{Langseth2003}。在因果学习框架下将两种来源的数据结合，得到可能的因果网络模型的概率分布。根据这一概率分布评估变量排序的不确定性，通过专家知识获取来收集新的数据提供重要证据，从而最大限度地降低估计变量排序的不确定性。

通过限制网络中节点的父集（入度）来缩小搜索空间的大小，证明了一般父集合很大的网络往往得分会比较低，在实际应用中很少有这样的网络结构并且很少有足够的数据来支撑\cite{Friedman2013}。另有研究使用了两种约束，一个节点的入度和节点之间的弧度，该方法可以极大地降低算法的时间和内存开销\cite{DeCampos2009}。部分研究将基因表达数据学习因果网络的各种结构限制编码为先验知识，通过描述一些类型的结构限制来解决额外先验知识的使用，考虑了三种类型的结构约束，包括存在性约束、缺失限制以及排序结构限制，证明了研究从表达数据中针对因果网络结构学习算法能够在利用结构约束的情况下利用这种先验知识获得更好结果\cite{De2019}。一些学者将领域知识的结构限制区分为硬约束和软约束，一些学者通过有向边的相关约束减少搜索空间，这些所提出的约束都属于硬约束，强制模型必须满足这些约束限制条件\cite{Wu2015}。另有研究提出了软约束，通过计算专家指定结构与候选结构之间编辑距离的先验分布来计算网络结构的相对后验概率，选取最优的候选结构\cite{Heckerman1995}。还有些学者通过将领域知识内的特定约束转换为规则来限制搜索策略的进行，利用特定的领域知识表示为规则，利用规则来约束因果结构网络的构建，减少网络的搜索空间\cite{赵建喆2020}。

\subsection{研究现状分析与总结}
综上所述，当前因果发现领域已在基于观测数据的结构学习及先验知识融合方面取得了丰硕成果，特别是基于连续优化的框架有效突破了高维场景下的组合搜索瓶颈，而知识引导策略则为数据匮乏场景提供了新的解决思路。然而，纵观现有研究，仍存在以下亟待突破的局限性：

首先，因果表示能力不足，重结构轻机制。现有基于连续优化的方法大多局限于学习有向无环图，仅能回答变量间“是否存在因果关联”，却无法揭示“如何影响”的内在机理。这种对因果函数的忽视，导致模型难以支持精确的干预决策与反事实推理。尽管部分方法尝试引入神经网络拟合非线性关系，但多层感知机的黑箱特性使得学到的因果机制缺乏数学透明性，无法提供人类可读的符号表达。

其次，知识融合深度不够，缺乏自适应引导机制。现有的知识融合方法多将领域知识视为硬性的结构约束或初始化先验，缺乏知识与数据驱动模型的深层交互。这种浅层融合不仅难以克服专家知识的主观性与不完整性，还限制了算法在复杂非参数场景下的泛化能力。特别是在面对真实世界中普遍存在的非参数或半参数数据时，现有方法难以利用有限参数模型精准显式地建模复杂的因果函数分布。

最后，求解器架构与因果可解释性需求不匹配。传统深度学习求解器基于通用逼近定理设计，其稠密连接与点乘操作隐式编码输入交互，天然契合统计相关性建模，却与因果机制的稀疏性、可分解性假设存在冲突。这导致即便在结构准确的情况下，模型仍难以解释变量间的具体作用形式（如线性、非线性、阈值效应等），限制了因果模型在医疗、政策等高可靠性要求场景中的落地应用。

因此，探索一种能够联合发现因果结构与函数、具备可解释性且能深度融合先验知识的因果发现新方法具有重要意义。本文旨在引入 Kolmogorov-Arnold 表示定理驱动下的神经网络，利用其天然的符号化潜力与函数分解特性，构建结构与函数联合发现的新方法，以解决上述问题。


\section{本文主要研究内容及创新点}
\subsection{主要研究内容}
上一节内容系统的介绍了基于连续优化的因果发现方法和基于先验知识融合的因果发现方法。本文以两种因果发现方法为研究对象，着力于填补其中的研究空白。研究内容聚焦两大挑战：1、现有因果发现方法的研究结果通常用一个有向无环图来表示，缺乏对因果函数的刻画，存在因果关系表示能力不足的局限。原因在于黑盒神经网络与因果发现的融合方法缺乏对因果关系的显式建模，因果函数被表示为一个神经网络，不透明且无法解释，这种表示方法忽视了结构因果模型中因果函数的重要性；2、面对真实场景中普遍存在的非参数或半参数因果数据，有限参数的因果函数模型难以精准刻画复杂数据分布，且缺乏有效融合先验知识的机制。图\ref{本文的研究内容及相互关系}展示了本文的研究内容及其相互关系。其中，针对挑战1，本文提出了一种基于KAN的连续优化算法，该算法不仅实现了因果图的端到端发现，而且将因果函数映射通过Kolmogorov-Arnold表示定理来显式建模；针对挑战2，本文提出了一种基于KAN的知识融合算法，该算法实现了深度KAN来显式建模非参数和半参数数据的因果函数，使得在真实场景下具有实用性和可解释性。下面详细介绍两个研究内容。
\begin{figure}[H]
	\centering
	\includegraphics[width=\linewidth]{研究内容及相互关系}
	\caption{本文的研究内容及相互关系}
	\label{本文的研究内容及相互关系}
\end{figure}

{\hei 1.基于KAN的连续优化因果发现算法研究}

针对现有连续优化方法仅能识别因果结构存在性而无法揭示变量间具体作用规律的问题，本研究提出一种基于 KAN 的连续优化因果发现方法。传统方法通常采用多层感知机（MLP）作为求解器，其黑箱特性导致学到的因果函数不透明，难以支持后续的干预决策。本研究利用 KAN 基于 Kolmogorov-Arnold 表示定理的可解释性优势，构建轻量级 KAN 架构作为因果发现求解器。该研究内容核心在于实现因果图结构与因果函数映射的端到端联合学习：一方面，通过 KAN 边的稀疏性约束学习因果邻接矩阵，确定变量间的依赖结构；另一方面，利用 KAN 可学习的 B 样条基函数显式拟合变量间的非线性因果机制。此外，针对加性噪声模型在高维复杂场景下结构学习精度下降的问题，本文分析了现有算法在节点独立性矩阵刻画上的不足，提出了适配 KAN 结构的节点独立性矩阵表示方法及稀疏化正则策略。实验表明，该算法在加性模型数据上的结构学习精度优于现有主流方法。具体研究细节详见本文第三章。

{\hei 2.基于KAN的知识融合因果发现算法研究}

针对真实世界数据往往符合非参数或半参数分布特点，且单纯数据驱动方法在数据稀缺或噪声较大时鲁棒性不足的问题，本研究进一步提出基于 KAN 的知识融合因果发现算法。该研究内容旨在突破有限参数模型对复杂因果函数的表达限制，并探索先验知识在因果发现中的深层融合机制。研究提出将网络架构解耦为因果知识嵌入层与函数拟合层：因果知识嵌入层通过约束第一层激活值，将已知的因果父节点关系编码为二进制掩码，强制网络反映$PA_i \rightarrow Y_i$的映射约束，使模型成为结构因果模型的透明估计器；函数拟合层则采用深度 KAN 架构，在知识嵌入层阻断非因果变量信息传导的基础上，释放深层网络的拟合能力以逼近复杂的非参数因果函数。在 Sachs 真实数据集上的实验验证表明，嵌入因果知识不仅实现了对复杂因果函数的显式建模，还作为一种结构正则化手段有效防止了过拟合，揭示了算法在现实场景中的应用价值。具体研究细节详见本文第四章。

\subsection{主要创新点}

本文围绕“因果函数显式建模”这一核心命题，提出两项具有递进关系的创新工作。二者均以 Kolmogorov-Arnold 网络为理论基石，突破了传统因果发现止步于图结构识别的局限，将研究范式从“因果存在性判断”系统性推进至“因果作用机制解析”。

{\hei 1.因果结构与函数的联合发现：面向加性模型的轻量级KAN优化框架}

提出了 DAG-KAN 方法，一种基于 Kolmogorov-Arnold 网络的因果结构与函数联合学习框架。其核心创新在于突破了传统连续优化方法仅能输出 DAG 拓扑的局限，首次实现了因果结构与因果作用规律（函数映射）的端到端联合学习。不同于依赖黑箱 MLP 的传统方法，DAG-KAN 采用轻量级单层 KAN 架构，通过可学习的 B 样条基函数显式建模变量间的非线性交互，使因果机制具备数学透明性。该方法特别适用于加性因果模型，在高维复杂场景下相比现有最优算法（如 CAM）结构汉明距离（SHD）平均降低 22.6\%以上，且无需依赖预处理或后剪枝步骤。DAG-KAN 将因果发现从结构识别推进到机制解析，为需要理解“如何干预”的领域提供了可操作的因果知识。

{\hei 2.因果知识融合的函数机制解析：面向非参数和半参数数据的深层KAN正则化架构}

提出了 Causal-KAN 方法，一种将先验因果知识嵌入神经网络架构的因果函数学习框架。其核心创新在于改造KAN网络结构，通过在输入层引入“因果掩码”机制，将部分已知的因果父节点关系硬编码为网络连接规则，仅允许真正的因果父变量激活对应的样条函数，非因果变量被物理隔离。这种设计构建了“知识嵌入层 - 函数拟合层”的解耦架构，在保障因果忠实性的同时释放了深层网络的函数逼近能力，使其能够处理非参数和半参数数据。实验表明，因果约束本身作为一种结构正则化有效防止了过拟合，且学到的样条函数可直接可视化为人类可读的因果作用曲线。该方法将因果发现从“关系存在性判断”提升至“作用机制解析”，特别适用于精准医疗、政策干预等需要理解具体影响机制的场景。

\section{本文的组织结构}
本文的结构安排如下：

第一章为绪论。

第二章为背景知识与理论基础概述。

第三章为基于KAN的连续优化因果发现算法研究。

第四章为基于KAN的知识融合因果发现算法研究。

第五章为总结与展望。
